\scanpage{73}
\mspage{17}

\uline{Lemme} \inserted{On a} $\sigma_X^{-1}(\{x\})\simeq{}_{x\backslash}X,\quad \beta_X^{-1}(\{y\})\simeq X_{/x}$. \deleted{\unsure{Le}}

Le foncteur $\sigma_X$ est \inserted{Cat-}fibrant à fibres \deleted{\unsure{asphé}} contractibles (donc asphériques, donc $W$-asphériques), le foncteur $\beta_X$ cofibrant à fibres contractibles (donc asphériques donc $W$-asphériques). Ce sont donc \inserted{aussi} des Hot-équivalences a fortiori des $W$-équivalences.
% slanted
\margin{En fait, ce sont même des \deleted{\unsure{équiv. d'}} \unsure{homotopismes}, cf. \unsure{ci-après|ci-dessous} (un argument plus intrinsèque \unsure{s'invite|s'insère} \unsure{naturellement dans|en particularisant} le cas de $X'=X$)}

\uline{Corollaire 1} Le foncteur $i_X:X\unsure{\text{---}|\to}\underline{\mathrm{Fl}}(X)$ est une Hot-équivalence (puisque \unsure{sa composée|son composé} avec la Hot-équivalence \inserted{$\sigma_X$} est l'identité).

Voici la construction générale de factorisation de $f:X'\longrightarrow X$, correspondant au choix précédent d'une ``catégorie des chemins'' $\unsure{\mathcal C|C}(X)=\underline{\mathrm{Fl}}(X)$
\diagram{\begin{tikzcd}[ampersand replacement=\&,row sep=small,column sep=normal]
\overline{X}'\arrow[dr]\arrow[drr,"\bar f"]\arrow[dd,"\pi_f" description,double] \&\&\\
\&\underline{\mathrm{Fl}}(X)\arrow[r,"\beta_X"']\arrow[d,"\sigma_X"]\&X\\
X'\arrow[uu,dashed,bend left=27,"i_f"]\arrow[r,"f"']\&\deleted{\unsure{X}}\unsure{X}\&
\arrow[from=1-1,to=3-2,phantom,"\mathrm{cart}" description]
\end{tikzcd}}
% slanted
\margin{\uline{NB} $\pi_f\cdot i_f=\mathrm{id}$, et $i_f\pi_f\longrightarrow\mathrm{id}$, donc $i_f,\pi_f$ sont des homotopismes, \deleted{\unsure{donc}} inverses l'un de l'autre dans $\mathrm{Cat}\,\mathrm{hot}$. Cela précise cor.~3 page 18.}
où $i_f$ est déduit de la section $i$ de $\underline{\mathrm{Fl}}(X)$ sur $X$. Ainsi, $\overline{X}'$ est la catégorie des couples
\[ (x',u)\mid x'\in\operatorname{Ob}X',\quad \underset{\substack{\text{\tikz[baseline=-.4ex]\node[rotate=-90,inner sep=0pt] {$\scriptstyle\in$};}\\\mathrm{Fl}(X)}}{u}:f(x')\longrightarrow y \]
le foncteur $\bar f$ est défini par
\[ \bar f(x',u)=\mathrm{but}(u)=y \]
et $i_f$ par
\[ i_f(x')=(x',\mathrm{id}_{f(x')}:f(x')\overset{\mathrm{id}}{\text{---}} f(x')) \]
donc on a bien $f(x')=\bar f\deleted{\unsure{\varphi}}\,i_f(x')$.

% archive sheet number: 72
